\chapter{Review of Related Literature}

The rapid advancement of global research and innovation has spurred the growth of Intellectual Property (IP) and Intellectual Property Rights (IPR). With this, there has been an avenue for research towards Intellectual Property Management Systems (IPMS) in businesses, universities, and research institutions. This literature review builds the context for the current state of IP and IPR in the Philippines, and further delves into the discussion of utilizing digital technologies for IP management, starting from a global perspective, narrowing down to the university level. This review aims to identify and acquire an understanding of related concepts, existing solutions, research gaps, as well as frameworks that would be essential for the development of an Intellectual Property Management System (IPMS) tailored to the Technology Transfer and Business Development Office (TTBDO) institution of the University of the Philippines Visayas (UPV).

\section{IP Management and Technology Transfer}

The Philippines recognizes the vital role of IPs in creativity, innovation, and the overall development of the nation, economically and socially; it recognizes the social function that IPs possess. In 1997, the State enacted Republic Act No. 8923, otherwise known as the ``Intellectual Property Code of the Philippines,'' which holds the power and authority to enforce all laws related to intellectual property protection. With it, the Intellectual Property Office of the Philippines (IPOPHL) was established as the national government agency responsible for administering and implementing the laws. It also serves as the central authority for the protection, regulation, and promotion of IPs.

Although the establishment of IPOPHL laid the groundwork for the system of IP protection, it lacks the function to extract commercial value from IPs. Bader et al. (2012, as cited in Agostini et al., 2017) stated that ``The mere accumulation of IP assets does not grant any value appropriation from innovation or any return on investment; to this purpose, IP management not only becomes a critical challenge for firms, but also a practice they cannot disregard if they want to reap the benefits of their IP.'' In fact, institutions that generate and protect intellectual property struggle with capitalization and extracting commercial value from them (Agostini et al., 2016, as cited in Agostini et al., 2017). Because of this challenge in the IP field for institutions and businesses, the IP protection perspective of their research departments changed from economic and legal to strategic and managerial, as strategic IP-driven decisions have a substantial impact on institutions and businesses (Di Minin and Faems, 2013, as cited in Agostini et al., 2017).

Recognizing this need, the Philippine government enacted Republic Act No. 10055, otherwise known as the Philippine Technology Transfer Act of 2009, with an objective that ``aims to promote and facilitate the transfer, dissemination, and effective use, management, and commercialization of intellectual property, technology and knowledge resulting from R\&D funded by the government for the benefit of national economy and taxpayers'' (Republic Act No. 10055, 2009, as stated in Department of Science and Technology [DOST], 2009). The act provides a framework for the ownership, management, and commercialization of IPs that are generated from government-funded research and development. This act also mandated all Research and Development Institutions (RDIs) of the country to establish their own Technology Licensing Offices (TLOs) and/or Technology Business Development Offices (TBDOs) that can have their own governing policies, in support of the laws and policies of the IPOPHL and the nation, regarding Intellectual Property Rights (IPR) management and Technology Transfer (Republic Act No. 10055, 2009).

In accordance with the Philippine Technology Transfer Act of 2009 (R.A. 10055), the University of the Philippines established the UP Technology Licensing Office, which was subsequently renamed the Technology Transfer and Business Development Office (TTBDO) in 2011. It is the university’s primary institution and mechanism for managing, protecting, and commercializing IPs generated from academic research, innovation, and creative works. Its primary functions include overseeing invention disclosures, facilitating intellectual property registrations and license agreements, collaborating with legal offices to enforce intellectual property rights, aiding in research collaborations, and advising researchers on commercialization plans (Technology Transfer and Business Development Office [TTBDO], n.d.).

\section{Digital Systems in IP Management}

In this modern age, a thriving industry in research and innovation has emerged, resulting in the rapid generation of intellectual property (IP) and intellectual property rights (IPR). Firms recognize the importance of IP and IPR by using IP-driven corporate strategies. Hence, there is a need to study the practice and research of IP management (Holgersson \& Santen, 2018). According to Kolodyazhnaya et al. (2020), the digitalization of the economic system environment has significantly altered the processes of IP management, protection, and commercialization. The paradigm shift of IP management has introduced new mechanisms of collaboration and technology transfer among universities, businesses, firms, and other organizations, emphasizing the need to adopt a system-based approach in IP management. Their study stated that, in terms of IP and IPR, the competitiveness of an institution in a socio-economic environment heavily depends on its ability to adapt to the evolving technological and digital environment. This integration of digital technologies into IP management has led to the study and specialization of IP management systems (IPMS).

Gargate \& Momoya (2018) explored the methods and designs for developing a structured model for an IPMS, arguing that systematic management bridges the gap between IP generation and commercialization. Jain and Sharma (2006) extended that notion through their study on an organizational perspective of IPMS. They stated that an IPMS is a mechanism that funnels IPs into strategic, revenue-generating assets that have social and economic value. Their study argued that an IPMS has a policy and managerial role, and identified five major functions as follows:
\begin{itemize}
   \item Generation of IP
   \item IP portfolio management
   \item IP valuation
   \item Competitive assessment
   \item Strategic decision-making
\end{itemize}

Moreover, they conducted a case study of an IPMS integrated into the Institute of Indian Technology, Bombay (IITB). The study found that the IPMS significantly contributed to the institute by increasing the rate of IP applications and technology transfers, which encouraged the academe to explore new research programs.

Krejcar et al. (2020) further discussed the importance and utilization of IPMS through their study on software solutions (SWS) for IP management by university Technology Transfer Offices (TTOs). The study involved a systematic literature review of 11 identified research articles related to SWS for IP management, interviews with the TTO heads of four different universities, and an evaluation of 10 existing commercially available SWS for IP management. The study found that technology transfer units of business institutions have different mechanisms from TTOs of universities; as such, commercial SWS are not optimal options for IP management for university TTOs. Commercial IPMSs are cost-ineffective and too complex for university-level operations; thus, the recommended IPMS for a university TTO would follow the institution’s workflow, be cost-efficient, and be customized for further specific university needs.

\section{IP Management Practices in Academic Institutions}

\subsection{UPV IP Management Process}
The UPV TTBDO follows a structured process for intellectual property application, as outlined in their Citizen Charter. The process involves systematic steps from initial submission to filing with the IPOPHL, with clearly defined timelines and responsible personnel.

\subsubsection{Official Process Framework}


Table~\ref{tab:ttbdo-process} summarizes the Citizen's Charter steps, actions, timelines, and responsible personnel. The entire process from initial submission to IPOPHL filing is designed to be completed within 19 days and 1 hour. The Intellectual Property Officer handles the majority of technical work, including disclosure review, IPR assessment, document preparation, and application finalization, while the Administrative Aide manages payment.

\begingroup
\setlength{\tabcolsep}{4pt}
\renewcommand{\arraystretch}{1.05}
\small
\begin{table}[H]
\centering
\caption[TTBDO IP Application Process (Citizen's Charter summary)]{
   TTBDO IP Application Process (Citizen's Charter summary). It shows the TTBDO's workflow starting from the submission of disclosure forms up until filing it with IPOPHL. For each step in the workflow, the idealized processing times and standardized fees are shown.
}
\label{tab:ttbdo-process}
\begin{tabular}{|p{4.5cm}|p{4.5cm}|p{2.5cm}|p{2.5cm}|}
\hline
\textbf{Client Step} & \textbf{Agency Action} & \textbf{Fees To Be Paid} & \textbf{Processing} \\
\hline
1. Submit Invention Disclosure Form and Commitment Letter
& 1. Receive, record, and review request and forward to staff-in-charge
& None 
& 30 min \\
\hline
& 1.1 Review completeness of submitted disclosure
& None
& 2 days \\
\hline
& 1.2 Assess and recommend applicable IPR to the client
& None 
& 1 day \\
\hline
2. Send confirmation on recommended IPR with appropriate Deed of Assignment
& 2. Prepare pertinent IP application documents and send to client for review and feedback
& None 
& 12 days \\
\hline
3. Provide feedback on drafted application
& Finalize application and submit to IPOPHL
& None 
& 3 days \\
\hline
& Facilitate payment of application fee
& None 
& 1 day\\
\hline
& 3.2 Send proof of filing to the client
& None 
& 30 min \\
\hline
\multicolumn{4}{|r|}{\textbf{Total standard time:} 19 days \& 1 hour}\vline\\
\hline
\end{tabular}
\end{table}
\endgroup


Figures \ref{fig:General_IP_Flow} and \ref{fig:Copyright_IP_Flow} illustrate the workflow of the IP application process, which shows the role of the Intellectual Property Officer in technical evaluation and document preparation, as well as the multi-stage review and feedback cycle between the IP Officer and Technology Generators before final submission to IPOPHL. The figures highlight the different actors involved in each stage and show where these actors perform actions or hold responsibilities. Figure \ref{fig:General_IP_Flow} also shows different stages where the different application processes have similarities, while Figure \ref{fig:Copyright_IP_Flow} is separated as it involves physical documents and wet signatures, branching into a different type of process flow. 

\begin{figure}[H]               
   \centering                   
   \includegraphics[scale=0.18]{flowchart_general.png}    
   \caption[TM, UM/ID, Patent Application Flow Chart]{TM, UM/ID, Patent Application Flow Chart. It shows the step-by-step process for Trademark, Patent, Utility Model, and Industrial Design rights application facilitated by the TTBDO. The flowchart visualizes different actors such as the technology generator, TTBDO, and IPOPHIL, highlighting stages where they interact and when documents are turned over.}
    \label{fig:General_IP_Flow}
\end{figure}

\begin{figure}[H]               
   \centering                   
   \includegraphics[scale=0.35]{flowchart_copyright.png}    
   \caption[Copyright Application Flow Chart]{Copyright Application Flow Chart. The flowchart for the Copyright application process. It can be seen in the chart that there are steps that involve physical documents and wet signatures.}
    \label{fig:Copyright_IP_Flow}
\end{figure}

\subsection{Implementation Tools and Current Practices}

Through interviews with TTBDO personnel, specifically the Intellectual Property Officer (IP Officer), the actual implementation of the IP management process reveals specifics not contained within the Citizen Charter. This involves reliance on various digital tools and manual procedures, which present challenges to the system.

\subsubsection{Application Submission}
The TTBDO solicits disclosure through text messages and email announcements given to the mother units (colleges or departments), with the expectation that these units will subsequently communicate the information to their faculty and staff constituents. Alternatively, technology generators can schedule appointments through a Google Form\footnote{\url{https://docs.google.com/forms/d/e/1FAIpQLSe5-3752JG4qOiUdv93xWs1kdtW_ojH5uRqpAJTJvs7MtVfjQ/viewform}} available on the TTBDO website. Both approaches rely on indirect access which results in lower visibility and accessibility compared to direct submission portals.

Once technology generators have applied for IP, the exchange of IP-related documents between TTBDO and applicants is through Google Docs. While this allows for collaborative editing and reviewing of documents, it lacks formal workflow management. It is also an inefficient process since the two parties still need to manually communicate with each other, such as when a document is completed and ready for review.

\subsubsection{Tracking and Monitoring}
Application status tracking is managed manually through Google Sheets, where staff input updates at various stages of the process. The IP Officer noted that this manual data entry approach is vulnerable to input errors, which could result in data loss or inaccurate records, particularly since multiple users can access the same spreadsheet. While the IP officer is the only one to maintain the Google Sheets, some users can access the same spreadsheets since the email that is used is handled by multiple users.

Payments are also tracked in the system. Handling payments involves technology generators, TTBDO, and IPOPHL. TTBDO receives a statement of account from IPOPHL and then informs the Technology Generator of payment requirements. Either the technology generator or the TTBDO processes payment. Once payment is received by IPOPHL, IPOPHL will then send a receipt to TTBDO. This process requires careful coordination of these parties, yet lacks a centralized system to monitor payment status. This is one of the pain points identified by the IP Officer, wherein the lack of payment status causes potential delays or lost documentation.

\subsubsection{Dashboard and Portfolio Monitoring}
Portfolio monitoring and dashboard are also conducted through Google Sheets, where staff use formulas for data visualizations. Like in the problem with regards to tracking, the IP Officer emphasized that the dependency on sensitive cell formulas and manual inputs introduces risks of accidental data corruption or loss.

\subsubsection{Data Storage Infrastructure}
The current IP management system employs a fragmented approach for storage. Staff stores data in local machines, which are applications that are not yet processed, while those that are processed are stored in Google Sheets and Google Drive. Google Drive also serves as document storage and sharing for various IP-related documents. The IP Officer explained that due to the distributed nature of this storage infrastructure, they have adopted practices to maintain data integrity. This includes regularly downloading local backup copies of the entire repository.

\subsubsection{Alert and Notification System}
The IP Officer emphasized the complete absence of automated alerts or notification mechanisms as a critical gap in their current system. IPOPHL enforces strict two-month deadlines for responding to certain application stages. Without automated reminders or deadline tracking, staff must manually monitor calendars, set personal reminders, and cascade this information to the applicants if needed. This creates significant risks of missing critical deadlines, which could result in dropping of IP applications.

\subsubsection{Processing Time}
While the Citizen Charter specifies a standard timeline of 19 days and 1 hour, the IP Officer noted that the actual processing time could be longer in real life. The primary factor of such delay is the concentration of responsibilities on a single IP officer.

\section{IP Management in Other Universities}
Several universities, both local and international, have implemented varying approaches to IP management. Among UP constituents, a moderate digital divide exists between UPV TTBDO and other campuses, primarily at the initial disclosure submission stage. UP Diliman provides both downloadable and directly submittable Invention Disclosure Forms through their website\footnote{\url{https://ttbdo.up.edu.ph/forms/}}, allowing technology generators to complete and submit forms online. However, subsequent processing involves email communication or physical submission at their office \cite[p.~6]{upd2020ovcrd}. Similarly, UP Mindanao offers direct online submission through Google Forms accessible on their resources page\footnote{\url{https://upmin-ttbdo.up.edu.ph/resources}}, with subsequent processing following manual procedures \cite[pp.~70--75]{upmindanao2024citizenscharter}.

In contrast, UPV TTBDO requires either scheduling appointments before disclosure submission through Google Form or relies on email announcements cascaded through mother units. Both methods introduce additional steps compared to the direct submission portals of UP Diliman and Mindanao. However, beyond the initial submission stage, all three campuses process the application manually through email or face-to-face coordination. However, specific tracking and monitoring systems vary and are not fully documented in publicly available materials.

Other local universities, such as De La Salle University (DLSU) and Polytechnic University of the Philippines (PUP), continue to rely on manual procedures. DLSU requires physical submission of IP disclosure forms (De La Salle University, n.d.), while PUP only gives out a Google Drive link to access forms, suggesting the need to submit and process applications physically. While an in-person process ensures accountability and document integrity, it might introduce challenges, including limited accessibility, increased processing time, and redundancy in document storage.

In contrast, international universities such as Massachusetts Institute of Technology (MIT), Stanford, and Oxford University employ a fully digitalized system. The full lifecycle of an IP application is supported through an online portal, from disclosure submission to registration  (MIT TLO, 2024; OUI, 2013, p. 11). Their systems include features such as online submission portals, status tracking, and structured review workflows. 

For instance, the MIT Technology Licensing Office (TLO) outlines a digital process for submitting invention disclosures, tracking patent evaluations, and coordinating with inventors and licensees through a centralized management system (MIT TLO, Inventors’ Guide, 2021, pp. 4–7; Patenting and Commercializing Your Invention, 2020). Similarly, Stanford University’s Office of Technology Licensing describes an online disclosure submission and review workflow that routes cases through assigned licensing associates, ensuring traceable evaluation and communication (Stanford OTL, OTL’s Process, 2021). The University of Oxford’s Innovation Office also details an integrated digital platform with defined user roles, patent status monitoring, and reporting dashboards for IP portfolio management (Oxford Innovation, IP, Patents, and Licences Researcher Booklet, 2013, p. 11).

Table \ref{tab:comparative} highlights clear distinctions between TTBDO’s current IP management process and those implemented by other universities. While international institutions such as MIT, Stanford, and Oxford have developed fully digital, centralized, and automated IP management systems, most local universities still depend on manual processes, fragmented storage, and limited tracking mechanisms. This analysis emphasizes the need for a web-based IP management system that integrates online submission, guided workflows, automated tracking, centralized data storage, notification mechanisms, document management, role-based access, and reporting support.

\begin{landscape}
\begingroup
\setlength{\tabcolsep}{4pt}
\renewcommand{\arraystretch}{1.05}
\small
\begin{table}
\centering
\caption[Comparative Summary of IP Management Systems]{Comparative Summary of IP Management Systems. The table compares key features of IP management systems used in both local and international university-level institutions. The UPV TTBDO's current IP management setup is compared with these to identify its gaps and weaknesses.}
\label{tab:comparative}
\begin{tabularx}{\linewidth}{|X|X|X|X|X|}
\hline
\textbf{Feature} & \textbf{UPV TTBDO} & \textbf{Local Universities} & \textbf{International Universities} & \textbf{Gap/Insight} \\
\hline
Disclosure Form Submission
& Google Form (for scheduling) or email cascade through units
& UP Diliman: Direct online submission via website; UP Mindanao: Direct Google Forms; DLSU and PUP: Physical
& Integrated online portals
& UPV requires intermediary steps (scheduling or unit communication) vs. direct disclosure submission \\
\hline
Subsequent Application Processing
& Manual via email and Google Docs
& UP Diliman: Email or physical; UP Mindanao: Manual per Citizen's Charter; DLSU and PUP: Physical
& Integrated online portals with workflow management
& Need for digital workflow after initial disclosure \\
\hline
Tracking
& Manual via Google Sheets
& Methods not publicly documented
& Automated digital tracking systems
& Need for automated tracking and process monitoring \\
\hline
Data Management
& Fragmented (local and Google Drive)
& Not specified in public documentation
& Centralized databases managed by IP offices
& Need for integrated database to ensure data integrity \\
\hline
Notifications \& Alerts
& None
& None mentioned
& Structured electronic correspondence and system notifications
& Need automated notification system \\
\hline
Staffing Structure
& Single IP Officer
& Office staff
& Dedicated licensing officers and technology transfer managers
& Need for role-based access and distributed management \\
\hline
\end{tabularx}
\end{table}
\endgroup
\end{landscape}



\section{Proposed Digital IP Management System}

Based on the gaps identified in the current TTBDO process and the digital practices observed in other universities, this study proposed IRIS as a web-based IP management system for UPV TTBDO. Since this chapter served as the literature review and conceptual basis of the study, the features described in this section are framed as proposed system capabilities rather than as completed implementation results. The proposed system was intended to support TTBDO operations by improving submission, monitoring, document handling, communication, accountability, and access control within the IP application process.

The proposed system aimed to provide the following core capabilities:
\begin{enumerate}
   \item \textbf{Centralized Data Management and Storage:} IRIS was intended to consolidate application records, user information, status histories, attachments, and supporting documents into a unified database and storage environment. This addressed the fragmentation of data across local machines, Google Sheets, Google Drive, and email threads.
   \item \textbf{Guided Application Submission and Review:} The system was proposed to provide Technology Generators and TTBDO Admins with a guided workflow for starting IP applications, selecting the appropriate IP type, entering application details, identifying collaborators, uploading required files, and submitting records for review. This feature was intended to reduce reliance on email cascades, informal document exchange, and manual coordination.
   \item \textbf{Application Tracking and Status Monitoring:} IRIS was designed to support role-based application tracking through status records, registry views, and application-specific pages. This capability addressed the limitations of manually maintained spreadsheets and supported clearer monitoring of each application’s progress.
   \item \textbf{Notifications and Deadline Reminders:} The proposed system included in-app and email notification mechanisms to inform users of important application updates, collaborator-related actions, and approaching deadlines. These reminders were intended to help TTBDO monitor time-sensitive IPOPHL requirements and support adherence to the Citizen Charter timeline.
   \item \textbf{Document and Attachment Management:} The system was proposed to manage uploaded files, application attachments, and reference documents in a more structured manner. This included support for controlled access, file traceability, and organized storage of application-related documents.
   \item \textbf{Collaborator Management:} Since IP applications may involve multiple Technology Generators, IRIS was intended to support collaborator identification, linking, reporting, and removal. This feature addressed the need to represent shared ownership and participation in IP applications.
   \item \textbf{Dashboards and Reporting:} The proposed system included dashboards and reporting tools that would provide visibility into application counts, status distributions, filing trends, and other IP portfolio indicators. These features were intended to support both operational monitoring and strategic decision-making.
   \item \textbf{User Management and Role-Based Access Control:} IRIS was designed to use secure authentication and role-based access permissions for Technology Generators, TTBDO Admins, UP Officials, and guests. This ensured that sensitive application information would only be accessible to authorized users while still allowing appropriate public access to selected information.
   \item \textbf{Public Access Features:} The proposed system included a public-facing component where guests could access selected application information, application documents, and guide materials. This supported transparency and improved accessibility for users who needed general information about the IP application process.
   \item \textbf{Audit Trail and Accountability:} The system was intended to record relevant user and system actions through audit logs. This supported traceability, accountability, and administrative oversight in the management of IP records.
   \item \textbf{Controlled Interoperability:} IRIS was also proposed to support controlled interoperability with external systems through secured application programming interfaces. This capability reflected the need for future integration while maintaining access control and data protection.
\end{enumerate}

Through these proposed capabilities, IRIS was positioned as a customized IPMS for UPV TTBDO rather than a generic commercial IP management tool. The system scope was aligned with the specific operational gaps identified in the current process, particularly fragmented storage, manual tracking, limited notifications, indirect submission channels, insufficient public access, and the need for role-based governance. The detailed functional requirements, system architecture, and technical specifications are presented in Chapter 3.

\section[Rationale for Digital IP Management]{Rationale for Digital IP Management}

\subsection{Current System’s Challenges}

The concentration of responsibilities on a single IP officer limits operational capacity. IP office managers need to carefully study the landscape of emerging technologies and understand how these new technologies affect them as service providers, developing strategies to absorb and take advantage of new digital technologies, and to adapt operations to fully online services (WIPO IPOBSD, n.d.). Digitalization provides productivity improvements, cost reductions, and innovations (Kraus et al., 2021), while digital systems provide the scalability necessary to handle increased IP activity and adapt to evolving technological landscapes.

Building on this concern, the current system’s reliance on manual data entry through Google Sheets and fragmented storage across local machines and Google Drive creates significant vulnerabilities. Digital transformation enables organizations to move from manual or procedural business processes to fully online, responsive and modern processes that enhance productivity and service delivery (WIPO IPOBSD, n.d.). Centralized digital systems eliminate the risks of data corruption from sensitive cell formulas and multi-user access issues, while providing systematic data protection without needing to create individual backup practices.

In addition to data management issues, the current IPs’ lack of automated alerts for IPOPHL's strict two-month deadlines represents a critical operational risk. Automated notification systems are needed to ensure tracking of deadlines and informing applicants regarding any changes to their applications. This, in turn, mitigates the risk of dropped applications due to missed deadlines.

Furthermore, challenges also exist in the coordination of application documents, collaborators, and status-related communication. IP applications may involve multiple Technology Generators, TTBDO personnel, and external offices, but the current process relies heavily on email, shared documents, and manual follow-ups. Efficient integration can only be achieved through digital processes and collaborative tools (White, 2012, as cited in Kraus et al., 2021). A centralized system for application records, document exchange, status monitoring, and collaborator management can reduce delays and prevent loss of information across the workflow.

The combination of these issues highlights the need for a holistic transformation that addresses operational bottlenecks, improves transparency, and supports accountability across IP processes.

\subsection{Benchmarking Digital IP Practices}
While local universities maintain partially digital or manual systems, international institutions demonstrate the full potential of digital IP management. The comparative analysis revealed that MIT, Stanford, and Oxford utilize integrated platforms supporting the complete lifecycle of IP applications. Digital technologies not only impact the transformation of products, business processes, or sales, but entire business models as well (Hess et al., 2016, as cited in Kraus et al., 2021). Thus, implementing a comprehensive digital solution positions UPV alongside these leading institutions.

Beyond institutional benchmarking, broader economic implications further strengthen the rationale for digital transformation. Digital Intellectual Property (IP) Rights listed robust protection of IP as one of the fundamental factors that contribute to a country's attractiveness from an investment perspective (DCO, 2024). Enhanced digital IP management capabilities strengthen institutional competitiveness and attractiveness for research collaboration and investment partnerships.

By aligning local systems with international standards, UPV not only modernizes its IP management process but also strengthens its role as a driver of innovation within the national research ecosystem.

\subsection{Strategic IP Management}
Having a dashboard to replace the manual visualizations produced by formulas in Google Sheets enables not just a truthful data representation but also offers more ways to do data visualizations and thus data analytics. WIPO IPOBSD (n.d.) states that data analytics may be used to understand trends in the IP system, moving beyond traditional statistics to deeper analysis techniques that may be used to understand micro trends, such as the distribution of filings from SMEs, or evolving technologies that may become significant in the future.

This information can be used to improve services to users, enabling IP offices to better understand the behavior of diverse users and develop strategies to better support them. Also, moving away from manual data entry and monitoring tasks enables staff to focus on strategic portfolio analysis and improved service delivery.

By equipping IP offices with tools for insight generation and trend monitoring, digital systems transform management from administrative to strategic, contributing to long-term organizational learning and decision-making.

\subsection{Institutional Impact}
The lack of a digital solution has created significant barriers, increasing the time and effort required to navigate the system, further discouraging researchers from pursuing intellectual property protection for their innovation. This cumbersome and inefficient process not only hinders innovation but also discourages collaboration and progress in the fields of research and development (Man et al., 2025).

In this context, digital technologies serve as a new channel for information or interaction with customers and other stakeholders (Li, 2020, as cited in Kraus et al., 2021). A user-friendly digital platform reduces these barriers, promoting increased participation in IP protection and supporting research commercialization efforts.

Moreover, IT facilitates information sharing and coordination, enabling the creation of value networks (Katsamakas, 2014, as cited in Kraus et al., 2021). The reshaping of customer value propositions can be achieved by using information and data analytics to enhance, extend, or redefine the customer experience (Berman, 2012, as cited in Kraus et al., 2021).

As a result, automated workflow tracking, notifications, and reporting tools can help TTBDO monitor progress toward the Citizen Charter's 19-day timeline, while integrated tracking and reporting capabilities enable performance measurement and continuous improvement. If the transformation is managed effectively, IP offices will be key players in the next phase of economic development (WIPO IPOBSD, n.d.), serving both as enablers of value creation through IP and as exemplars of digital transformation for institutional services.

\section{Summary}
This section highlights the growing importance of IP management in universities, particularly under the mandates of the Philippine Technology Transfer Act of 2009. While the UPV TTBDO performs essential roles in IP protection and commercialization, its processes remain largely manual. TTBDO relies on fragmented tools such as Google Sheets, Drive, and email communication. Studies show that transforming these manual processes to digital systems enhances efficiency, data integrity, and collaboration in IP management, as evidenced by the fully digital systems of leading international institutions like MIT, Stanford, and Oxford. Thus, IRIS, a web-based IPMS tailored to UPV TTBDO, was proposed to streamline application submission and tracking, consolidate data and document management, support collaborator management, provide notifications and deadline reminders, enable role-based dashboards and reporting, maintain audit trails, provide selected public access, and allow controlled interoperability with external systems. In this way, the proposed IPMS was intended not only to reduce administrative inefficiencies but also to support TTBDO’s broader role in innovation management and technology transfer.

